\documentclass[12pt,a4paper]{article}
\usepackage[T2A,T1]{fontenc}
\usepackage[utf8]{inputenc}
\usepackage[english,russian]{babel}
\usepackage{XCharter}
\usepackage{amsmath,amssymb,mathtools}
\usepackage{array,longtable,booktabs}
\usepackage[margin=24mm]{geometry}
\usepackage[hidelinks]{hyperref}
\urlstyle{same}
\setlength{\parindent}{0pt}
\setlength{\parskip}{0.45em}
\setlength{\LTleft}{0pt}
\setlength{\LTright}{0pt}
\emergencystretch=2em
\setlength{\tabcolsep}{4pt}
\newcolumntype{P}[1]{>{\raggedright\arraybackslash}p{#1}}
\title{Идеи Н. П. Брусенцова: краткое саммари}
\author{}
\date{}
\begin{document}
\maketitle

\emph{Для читателей без подготовки по логике и для специалистов. Изложение отделяет тезисы Н. П. Брусенцова от поясняющих реконструкций и оценки исследователей.}

\section{В двух абзацах}

В 2010 году Н. П. Брусенцов различает необходимое следование \(x\Rightarrow y\) и материальную импликацию \(x\to y\), а равнозначность между ними ограничивает универсумом \(\text{УА}\) с постулатом сосуществования противоположностей. Это позиция автора, а не открытый вопрос, который он только предлагает исследовать. Содержательность отношения он объясняет тем, что качество \(y\), присущее \(x\)-вещам, необходимо содержится в их сущности; это авторское философское основание, а не независимое доказательство причинной или физической связи. В 2008 году он полемически называет материальную импликацию «типичной химерой» (PDF с. 4; соответствие этой страницы печатной пагинации не установлено). Сравнение, в котором ячейка \(x'y\) получает то же значение, что другие истинные комбинации в двузначной оценке \(x\to y\equiv x'\lor y\), — наш пересчёт таблицы, а не формулировка Н. П. Брусенцова. Точные прямые формулировки и страницы первоисточников приведены в основном конспекте и в чек-листе ревизии. \href{https://ternarycomp.cs.msu.ru/Papers/Alrek.pdf}{«Алгебраическая реконструкция силлогистики»}, PDF с. 1, 4; \href{https://ternarycomp.cs.msu.ru/Papers/brusentsov_sborn2010_2.pdf}{«Воссоздание аристотелевой безукоризненной силлогистики»}, печатные с. 88–89.

Его предложение — представить \(x\Rightarrow y\), отношение необходимого содержательного следования, через статусы четырёх сочетаний терминов и постулат сосуществования термина и его противоположности. В главе 2005 года профиль \(+-0+\) — алгебраическое представление отношения. В авторском универсуме переход к экстенсиональному чтению дан явно: классы \(xy\) и \(x'y'\) непусты, класс \(xy'\) пуст, а статус \(x'y\) привходящий. Сам знак \(+\) не является квантором; применение этого смысла к конечному списку людей — наш пример авторской трактовки. После устранения среднего термина автор получает силлогистические следствия в этой модели. Это результат при выбранных определениях и предпосылках, а не доказательство несостоятельности классической логики или исторической тождественности модели логике Аристотеля. \href{https://ternarycomp.cs.msu.ru/Papers/Rea.pdf}{«Реанимация аристотелевой силлогистики»}, PDF с. 1–3; печатные с. 140–142.

Здесь \(x,y\) — термины/признаки, \(x',y'\) — противоположные термины, а \(xy,xy',x'y,x'y'\) — сочетания в этом порядке. В двузначном примере штрих можно пояснять как отсутствие признака; в общей нотации он не тождествен автоматически обычному отрицанию. \(P,Q\) обозначают целые высказывания, поэтому \(P\to Q\), семантическое следование \(\Gamma\models Q\), выводимость \(\Gamma\vdash Q\) и отношение \(x\Rightarrow y\) принадлежат разным уровням. Здесь \(\Gamma\) — множество посылок. Семантическое следование говорит о сохранении истинности во всех моделях, удовлетворяющих посылкам; выводимость — о наличии доказательства по правилам конкретного исчисления. Их совпадение требует согласованной семантики и правил выбранной системы и не устанавливает истинность посылок. В статье 2010 года \(Axy\) — общеутвердительное суждение «Все \(x\) суть \(y\)»; \(A(x,y)\) в статье 2003 года — статусная функция, и одинаковую букву нужно читать по контексту.

Стандартные знаки в формулах этого саммари: \(\forall\) — «для каждого», \(\exists\) — «существует», \(\land\) — «и», \(\lor\) — «или», \(\neg\) — «не»; \(\to\) связывает условное высказывание, когда его части — целые высказывания. В точечной записи \(X(t)\to Y(t)\), \(X,Y\) — классы носителей признаков \(x,y\), а \(t\) — отдельный предмет; знаки \(X(t)\), \(Y(t)\) утверждают, что предмет обладает соответствующим признаком. В формулах Н. П. Брусенцова те же знаки могут применяться к терминам и классовым выражениям; такой контекст отмечен отдельно. Для множеств \(U\) — область предметов, \(X,Y\subseteq U\) означает подмножества этой области; \(\cap\) — пересечение, \(\cup\) — объединение, \(\setminus\) — разность множеств (элементы первого множества, не входящие во второе).

В работах встречаются разные кодировки, близкие по общей идее, но не тождественные по значениям знаков. В статье 2003 года \(A(x,y)\in\{1,0,\sigma\}\): \(\in\) означает «принадлежит множеству», фигурные скобки перечисляют значения; \(1\) — необходимая присущность, \(0\) — антиприсущность, \(\sigma\) — привходящий статус, при котором \(y\) не необходимо присуще и не антиприсуще \(x\). Автор не исключает с необходимостью ни \(xy\)-, ни \(xy'\)-вещи; это не просто нехватка знания о фактическом составе класса. Условие \(0<\sigma<1\) задаёт положение \(\sigma\) между крайними значениями, но источник не фиксирует конкретное число и не уточняет, допускаются ли разные числовые реализации этого статуса. В главе 2005 года обобщённая форма \(xy\lor\sigma x'y\lor x'y'\) помечает привходящую ячейку \(x'y\) через \(\sigma\); знаковая шкала этой главы записывает соответствующий статус нулём в профиле \(+-0+\). Ноль 2005 года поэтому не равен антиприсущему нулю функции 2003 года. Профиль \(+-++\) — наш отдельный пересчёт точечной материальной импликации по четырём сочетаниям, а не профиль Н. П. Брусенцова; в этом сравнении \(+\) означает истину, \(-\) — ложь. Одинаковые знаки \(+\) и \(-\) в двух профилях не имеют автоматически одного смысла.

\section{Объяснение без специальной подготовки}

\subsection{О чём говорят буквы \(x\) и \(y\)}

Представим, что мы рассматриваем предметы и два их признака. \(x\) означает первый признак, \(y\) — второй; штрих означает противоположный признак. В этом бинарном примере противоположность читается как отсутствие признака; в общей нотации Н. П. Брусенцова штрих не следует автоматически сводить к обычному отрицанию. Тогда предмет может попасть в одну из четырёх групп:

\begin{longtable}{P{0.260\textwidth}P{0.300\textwidth}P{0.360\textwidth}}
\toprule
\textbf{Группа} & \textbf{Первый признак} & \textbf{Второй признак} \\
\midrule
\endfirsthead
\toprule
\textbf{Группа} & \textbf{Первый признак} & \textbf{Второй признак} \\
\midrule
\endhead
\(xy\) & есть \(x\) & есть \(y\) \\
\(xy'\) & есть \(x\) & нет \(y\) \\
\(x'y\) & нет \(x\) & есть \(y\) \\
\(x'y'\) & нет \(x\) & нет \(y\) \\
\bottomrule
\end{longtable}

Например, можно взять \(x=\) «носит очки», \(y=\) «студент». Это наглядный пример. В главе 2005 года Н. П. Брусенцов сам связывает свою нормальную форму с непустотой классов \(xy\) и \(x'y'\), пустотой \(xy'\) и привходящим статусом \(x'y\). Поэтому экзистенциальное чтение в его модели не является нашей догадкой. Нашей является конкретная иллюстрация конечным списком участников; знаки профиля при этом не следует принимать за сами кванторы существования. \href{https://ternarycomp.cs.msu.ru/Papers/Rea.pdf}{«Реанимация аристотелевой силлогистики»}, PDF с. 1–3.

\subsection{Что значит профиль \(+-0+\)}

В главе «Реанимация аристотелевой силлогистики» четыре позиции идут в порядке \(xy,xy',x'y,x'y'\). Профиль \(+-0+\) задаёт такой рисунок:

\begin{longtable}{P{0.290\textwidth}P{0.620\textwidth}}
\toprule
\textbf{Сочетание} & \textbf{Статус в профиле} \\
\midrule
\endfirsthead
\toprule
\textbf{Сочетание} & \textbf{Статус в профиле} \\
\midrule
\endhead
\(xy\) & необходимо включённое \\
\(xy'\) & необходимо исключённое \\
\(x'y\) & привходящее: может быть включено или исключено \\
\(x'y'\) & необходимо включённое \\
\bottomrule
\end{longtable}

Здесь знак \(0\) занимает \textbf{одну позицию} — сочетание \(x'y\). Сам по себе он не означает, что всё высказывание «если \(P\), то \(Q\)» получило третье значение. В этой главе Н. П. Брусенцов действительно называет промежуточный статус классового выражения третьим значением истинности; в его знаковой шкале необходимые включение, исключение и привходящий случай обозначаются соответственно \(+1,-1,0\). Вектор \(+-0+\) показывает статус четырёх комбинаций, а не таблицу значений двух целых высказываний.

В главе этот рисунок выражен обобщённой дизъюнктивной формой \(xy\lor\sigma x'y\lor x'y'\): знак \(\lor\) соединяет здесь выражения алгебры классов, клетка \(xy'\) необходимо исключена, а \(x'y\) помечена как привходящая. Здесь \(\sigma\) — авторское обозначение промежуточного статуса в этой форме; в знаковой шкале главы тот же тип статуса в соответствующей позиции записан нулём. Это не равенство числа \(\sigma\) нулю и не перевод в кодировку функции \(A(x,y)\) 2003 года, где ноль означает антиприсущность. Слова «необходимо включено», «необходимо исключено» и «привходящее» далее обозначают статусы в модели Н. П. Брусенцова, а не автоматически метафизическую, физическую или причинную необходимость. Формулу нельзя читать как таблицу истинности условного высказывания между двумя пропозициональными буквами.

В учебном примере с очками и студентами мы реализуем авторские условия непустоты и пустоты в конкретной конечной области. Такая иллюстрация полезна для проверки понимания, но не доказывает, что исходный аппарат Н. П. Брусенцова сводится к проверке списка людей.

\subsection{Зачем вводится «сосуществование противоположностей»}

Н. П. Брусенцов требует, чтобы в принятом универсуме сосуществовали каждый первичный термин и его противоположность. В нашем бинарном примере это означает, что есть предметы и с признаком, и без него. Это постулат сосуществования противоположностей: он исключает случаи, когда термин пуст или охватывает всю область. Поэтому правила о связи признаков зависят не только от записи формул, но и от предпосылок о том, какие предметы существуют.

Эта предпосылка существенна. Если класс \(x\) пуст, фраза «все \(x\) обладают \(y\)» в современной логике первого порядка считается истинной по пустому случаю. Но такая истинность сама по себе не говорит, что между признаками \(x\) и \(y\) есть содержательная связь. Для вывода от общего суждения к частному требуется отдельно указать, что \(x\)-предметы существуют.

\section{Что Н. П. Брусенцов предлагает логически}

\subsection{Не смешивать четыре разных уровня}

\begin{enumerate}
\item \textbf{Предметы и признаки.} Это то, о чём говорится: люди, числа, классы предметов и свойства.
\item \textbf{Термины и классы.} Буквы \(x,y\) обозначают термины/признаки; в формулах рассматриваются также классы предметов, которым эти признаки присущи. Символ термина и его класс не следует без оговорки объявлять одним и тем же объектом.
\item \textbf{Суждения.} Например, «всякий \(x\) есть \(y\)» — утверждение о классах.
\item \textbf{Вывод.} Правило показывает, когда из одной или нескольких посылок разрешено получить заключение.
\end{enumerate}
Н. П. Брусенцов в первую очередь задаёт собственное отношение \(x\Rightarrow y\) \textbf{между терминами и классами}. Его запись \(x\to y\) в алгебре терминов — двузначное материально-импликативное выражение для классовых сочетаний; она не является автоматически пропозициональной формулой \(P\to Q\). В то же время \(+-0+\) — не само отношение \(x\Rightarrow y\), а его статусное представление в главе 2005 года. Материальную импликацию \(P\to Q\) следует отличать и от семантического следования \(\Gamma\models Q\), которое требует сохранения истинности во всех допустимых моделях, и от выводимости \(\Gamma\vdash Q\), то есть наличия доказательства по правилам выбранной системы. Формальная корректность вывода при принятых посылках не гарантирует истинности посылок.

Материальная импликация может быть истинна, когда её предпосылка ложна. Если дополнительно представить \(X,Y\) как признаки отдельных предметов и развернуть условие \(X(t)\to Y(t)\) по четырём сочетаниям принадлежности, получится профиль \(+-++\): ложна только комбинация \(xy'\). Это наш пересчёт обычной импликации в предметной интерпретации, а не профиль, введённый Н. П. Брусенцовым. В его представлении \(x'y\) получает привходящий статус вместо безусловной истинности; при этом само \(x\Rightarrow y\) остаётся отношением, а не таблицей истинности. Сравнение показывает отличие двух кодировок, но само по себе не опровергает семантическое следование или правила вывода. Современные релевантные логики, модальные условные высказывания, силлогистика и теории объяснения исследуют разные части этой проблемы, но не составляют одну общепринятую «теорию содержательности» для всех контекстов.

\subsection{Почему в работах встречаются разные нули}

Обозначения нельзя переносить из одной публикации в другую без проверки:

\begin{longtable}{P{0.260\textwidth}P{0.300\textwidth}P{0.360\textwidth}}
\toprule
\textbf{Источник} & \textbf{Кодировка} & \textbf{Смысл} \\
\midrule
\endfirsthead
\toprule
\textbf{Источник} & \textbf{Кодировка} & \textbf{Смысл} \\
\midrule
\endhead
«Трехзначная интерпретация силлогистики Аристотеля» (2003) & \(A(x,y)\in\{1,0,\sigma\}\) & \(1\) — необходимая присущность; \(0\) — антиприсущность; \(\sigma\) — привходящий статус, причём \(0<\sigma<1\), но конкретное числовое значение не указано. \\
«Реанимация аристотелевой силлогистики» (2005) & \(+1,-1,0\) & Необходимое включение, необходимое исключение, привходящий статус. Автор описывает \(\sigma\) как промежуточное значение в обобщённой ДНФ, затем использует знаковую шкалу. \\
Профиль четырёх сочетаний (2005) & \(+-0+\) & Знаки приписаны позициям \(xy,xy',x'y,x'y'\); \(0\) в третьей позиции отмечает привходящий статус \(x'y\). \\
\bottomrule
\end{longtable}

Поэтому верно одновременно следующее: Н. П. Брусенцов использует третье значение в собственной алгебре классовых выражений; это \textbf{не} делает его таблицу таблицей истинности произвольных предложений \(P\to Q\). Кроме того, числовой \(0\) функции \(A(x,y)\) 2003 года нельзя автоматически отождествлять с \(0\) знаковой шкалы 2005 года: в этих записях числам назначены разные статусы.

\subsection{Как это связано с силлогизмом}

В аристотелевском силлогизме средний термин связывает две посылки. Н. П. Брусенцов стремится выразить такое связывание через распределение классов и затем алгебраически получить заключение после устранения среднего термина. Это не сводится к записи двух посылок через конъюнкцию и последующую материальную импликацию: структуру терминов и их классов нужно сохранить.

У Аристотеля есть и отдельный вопрос об объяснении: в демонстративном знании посылки должны не только влечь заключение, но и давать основание понимать, почему оно истинно. Валидность перехода и причинное объяснение факта — разные задачи. Сходство слов «необходимо» и «следует» не доказывает, что формализация Н. П. Брусенцова совпадает с аристотелевской.

\section{Что в этом подтверждается, а что остаётся спорным}

\subsection{Сильная сторона программы}

\begin{itemize}
\item Она напоминает, что правила рассуждения применяются к понятиям с содержанием и к предметной области, а не к буквам в вакууме.
\item Она делает явными предпосылки о существовании терминов и их противоположностей, которые иначе могут оставаться скрытыми.
\item Она различает необходимое включение, необходимое исключение и привходящий случай в модели классов.
\item Она связывает формальную логику с алгеброй и программной реализацией, которую можно проверять независимо от философских выводов автора.
\end{itemize}
\subsection{Что не следует считать доказанным}

\begin{itemize}
\item Что материальная импликация, семантическое следование или классическая логика в целом «не содержательны» во всех применениях.
\item Что введённая Н. П. Брусенцовым схема является единственным или общепринятым способом формализовать содержательную связь.
\item Что его интерпретация исторически тождественна логике Аристотеля. Это реконструкция автора; её нужно сравнивать с греческим текстом, контекстом и другими исследованиями.
\item Что троичная компьютерная техника сама по себе доказывает философскую или логическую правильность трёхзначной схемы. Реализация показывает, что операции можно кодировать и выполнять; обоснование смысла операций — отдельный вопрос.
\item Что троичные разряды вычислительной техники и статусы логической модели — один и тот же объект только потому, что оба используют три значения. Кодирование в машине и семантическая интерпретация кода требуют разных обоснований.
\end{itemize}
В статье 2003 года, её английской аннотации и более поздних работах встречаются разные числа правильных модусов: 192, 190 и 128. Пока основания подсчёта не приведены к одной системе обозначений и набору предпосылок, ни одно из этих чисел не следует объявлять окончательным.

\subsection{Аристотель и закон исключённого третьего}

В «Метафизике» Γ.7 Аристотель утверждает, что между противоречащими высказываниями нет промежуточного: о предмете нужно либо утверждать предикат, либо отрицать его. Это явная формулировка, близкая по содержанию к закону исключённого третьего; привычная символическая запись \(P\lor\neg P\) — более поздний язык. Поэтому тезис Н. П. Брусенцова, что закон несправедливо приписали Аристотелю как постулат Хрисиппа, нельзя принимать без уточнения: историческое происхождение термина и символической записи — не то же самое, что наличие у Аристотеля соответствующего принципа.

В «Об истолковании» 9 Аристотель обсуждает будущие случайные события: необходимо, что морское сражение завтра будет или не будет, но ни один исход не становится от этого необходимым сам по себе. Это ограничивает простой фаталистический вывод о будущих событиях; оно не превращает автоматически \(\sigma\) Н. П. Брусенцова в собственную логику Аристотеля.

\section{Для специалиста: проверяемая формальная программа}

Чтобы проверить предложение, нужно задать отдельно:

\begin{enumerate}
\item область предметов \(U\) и два термина как подмножества \(X,Y\subseteq U\);
\item четыре клетки \(X\cap Y\), \(X\setminus Y\), \(Y\setminus X\), \(U\setminus(X\cup Y)\);
\item функцию, приписывающую клеткам один из авторских статусов, и точную семантику каждого статуса;
\item постулат сосуществования и область, в которой он обязателен;
\item операции над статусами/шкалами, включая конъюнкцию, дизъюнкцию, инверсию и исключение среднего термина;
\item условия, при которых соответствующие операции воспроизводят конкретные силлогизмы.
\end{enumerate}
После этого можно доказать корректность отдельных правил, найти конечные контрмодели к неверному правилу и пересчитать модусы при явно заданных соглашениях. Такая проверка подтвердит следствия \textbf{из записанной модели}. Она не докажет, что модель правильно понимает Аристотеля или что выбранное определение следует называть «содержательным».

В отдельном заключении рассматривается, когда для этого полезен Lean и какие части дешевле сначала проверить обычным перебором.

Руководство \href{Как-читать-формулы-Брусенцова.tex}{«Как читать формулы Н. П. Брусенцова»} отдельно разбирает обозначения статьи 2010 года и отличает их от более явно пояснённой записи через интегральную дизъюнкцию в публикации 2001 года.

\section{Библиография к этому саммари}

\subsection{Работы Н. П. Брусенцова}

\begin{itemize}
\item Н. П. Брусенцов. \href{https://computer-museum.ru/histussr/trilog0.htm}{«Трехзначная интерпретация силлогистики Аристотеля» (2003)}. Архивная страница; \href{../sources/pdfs/Брусенцов_Трехзначная-интерпретация-силлогистики-Аристотеля_2003_версия-музея.pdf}{производная PDF-копия в корпусе}.
\item Н. П. Брусенцов. \href{https://ternarycomp.cs.msu.ru/Papers/Rea.pdf}{«Реанимация аристотелевой силлогистики» (2005), глава из \emph{Реставрации логики}, с. 140–145}. \href{../sources/pdfs/Брусенцов_Реанимация-аристотелевой-силлогистики_2005-глава-из-Реставрации-логики.pdf}{PDF в корпусе}.
\item Н. П. Брусенцов. \href{https://ternarycomp.cs.msu.ru/Papers/Soder_sled.pdf}{«Содержательное следование, сосуществование противоположностей и формальные импликации» (2006)}. \href{../sources/pdfs/Брусенцов_Содержательное-следование-и-формальные-импликации_2006.pdf}{PDF в корпусе}.
\item Н. П. Брусенцов. \href{https://ternarycomp.cs.msu.ru/Papers/Alrek.pdf}{«Алгебраическая реконструкция силлогистики» (2008)}. \href{../sources/pdfs/Брусенцов_Алгебраическая-реконструкция-силлогистики_2008.pdf}{PDF в корпусе}.
\item Н. П. Брусенцов. \href{https://ternarycomp.cs.msu.ru/Papers/brusentsov_sborn2010_1.pdf}{«Трехзначность отношения следования» (2010), печатные с. 86–87}. \href{../sources/pdfs/Брусенцов_Трехзначность-отношения-следования_2010.pdf}{Локальный PDF}.
\item Н. П. Брусенцов. \href{https://ternarycomp.cs.msu.ru/Papers/brusentsov_sborn2010_2.pdf}{«Воссоздание аристотелевой безукоризненной силлогистики» (2010), печатные с. 88–91}. \href{../sources/pdfs/Брусенцов_Воссоздание-аристотелевой-силлогистики_2010.pdf}{Локальный PDF}.
\item Н. П. Брусенцов. \href{https://ternarycomp.cs.msu.ru/Papers/Brusen_IMI14_Neobh_sled.pdf}{«Аристотелево необходимое следование» (2014)}. \href{../sources/pdfs/Брусенцов_Аристотелево-необходимое-следование-и-ДНФ_2014.pdf}{PDF в корпусе}.
\item Н. П. Брусенцов, Ю. С. Владимирова. \href{https://ternarycomp.cs.msu.ru/Papers/MMRO12.pdf}{«Троичная компьютеризация логики» (2005)}. \href{../sources/pdfs/Брусенцов_Владимирова_Троичная-компьютеризация-логики_2005_ММРО-12.pdf}{PDF в корпусе}.
\end{itemize}
\subsection{Аристотель и современные справочные работы}

\begin{itemize}
\item Аристотель. \href{https://classics.mit.edu/Aristotle/prior.1.i.html}{«Первая аналитика», I.1}.
\item Аристотель. \href{https://classics.mit.edu/Aristotle/posterior.1.i.html}{«Вторая аналитика», I.2}.
\item Аристотель. \href{https://aristotle.johnhboyer.com/Meta/book/4/}{«Метафизика», Γ.7}.
\item Аристотель. \href{https://www.st-andrews.ac.uk/~mnat/~ball0888/salamis/interpretatione.html}{«Об истолковании», гл. 9}.
\item Stanford Encyclopedia of Philosophy. \href{https://plato.stanford.edu/entries/aristotle-logic/}{Aristotle’s Logic}, \href{https://plato.stanford.edu/entries/logical-consequence/}{Logical Consequence}, \href{https://plato.stanford.edu/entries/logic-relevance/}{Relevance Logic}.
\end{itemize}
\section{Связанные документы проекта}

\begin{itemize}
\item \href{../research/Карта-работ-Брусенцова.tex}{Карта корпуса и доступности работ}.
\item \href{../research/Ревизия-корпуса-Брусенцова.tex}{Ревизия тезисов по первичным материалам}.
\item \href{../research/Проверка-идей-Брусенцова-в-Lean.tex}{Проверка формализаций в Lean}.
\item \href{../research/Полный-указатель-трудов-Брусенцова.md}{Полный указатель найденных работ}.
\end{itemize}
\end{document}
